\documentclass[11pt]{article}
\usepackage{amsmath,amssymb,amsthm}
\usepackage[margin=1.1in]{geometry}
\newtheorem{definition}{Definition}
\newtheorem{lemma}{Lemma}
\newtheorem{proposition}{Proposition}
\theoremstyle{remark}
\newtheorem{remark}{Remark}
\newtheorem{example}{Example}
\newcommand{\N}{\mathbb{N}}
\newcommand{\R}{\mathbb{R}}
\newcommand{\Z}{\mathbb{Z}}
\newcommand{\F}{\mathbb{F}}
\newcommand{\Aut}{\operatorname{Aut}}
\newcommand{\Sym}{\operatorname{Sym}}

\title{Crystals and their shadows:\\ definitions for a programme on the multiplication crystal}
\author{Carl Gribble\thanks{Programme Note 1, draft v0.1. Companion to the papers
\cite{CompanionI,CompanionII,CompanionIII} and the essay \cite{Essay} (same author), whose vocabulary
this note fixes. Developed in collaboration with Claude (Anthropic). Responsibility for the content
rests with the author. A full disclosure of AI use appears at the end of the note.}}
\date{August 2026 --- draft}

\begin{document}
\maketitle

\begin{abstract}
This note fixes the definitional frame for the multiplication-crystal programme: a \emph{crystal} is
a free commutative monoid equipped with an additive length functional whose axis spectrum is locally
finite. The definition is not new --- it is Knopfmacher's arithmetical semigroup written additively,
and Beurling's generalized number system in monoid form --- but writing it once lets the programme's
statements be made at full generality and specialized to its three standing instances: the integers,
the function fields, and the Beurling continuum between them, with the graph universe as a preview.
We define the invariant package the essay's chapters compute --- the zeta/partition function, the
layer decomposition and tower exponent, the Linnik prime detector, the shadow combs and their
diffraction, the automorphism group, and the lattice/non-lattice dichotomy --- prove the three
elementary structural facts the programme leans on (local finiteness of the shadow, the length
multiset as a complete isomorphism invariant, existence and uniqueness of the tower exponent), and
record which work package consumes which invariant. One observation falls out of the definitions and
deserves emphasis: measured by lengths alone, the integer crystal is already perfectly rigid, while
the function-field crystal retains an enormous symmetry group --- the symmetry asymmetry between the
two universes is visible before any arithmetic is welded on.
\end{abstract}

\section{The object}

\begin{definition}[Crystal]\label{def:crystal}
A \emph{crystal} is a pair $C = (M, \ell)$ where $M$ is the free commutative monoid on a nonempty,
at most countable set $P$ of generators (the \emph{axes} of $C$), and
$\ell \colon M \to [0, \infty)$ is a monoid homomorphism into $(\R_{\ge 0}, +)$ with $\ell(p) > 0$
for every $p \in P$, subject to the finiteness axiom
\[
\text{(F)}\qquad \#\{\, p \in P : \ell(p) \le x \,\} \;<\; \infty \quad\text{for every } x \ge 0 .
\]
We call $\ell$ the \emph{length functional}, $\ell(p)$ the \emph{axis lengths}, and
$|m| := e^{\ell(m)}$ the \emph{size} of $m$. The \emph{shadow} of $C$ is the image of $M$ under
$\ell$ counted with multiplicity: the measure $\nu_C = \sum_{m \in M} \delta_{\ell(m)}$ on
$[0, \infty)$.
\end{definition}

Freeness means every $m \in M$ is a finite product $\prod_p p^{a_p}$ with uniquely determined
exponents --- unique factorization is part of the definition, not a theorem. The axes are the
\emph{primes of the crystal}; we reserve the word ``prime'' for them.

\begin{lemma}[The shadow is locally finite]\label{lem:locfin}
Under \emph{(F)}, $N_C(x) := \#\{ m \in M : \ell(m) \le x \}$ is finite for every $x$.
\end{lemma}

\begin{proof}
If $p$ divides $m$ then $\ell(p) \le \ell(m)$, so an $m$ with $\ell(m) \le x$ uses only axes in the
finite set $P_x = \{p : \ell(p) \le x\}$. If $P_x$ is empty only $m = 1$ qualifies. Otherwise
$\delta = \min_{P_x} \ell > 0$, each exponent vector supported on $P_x$ has coordinate sum at most
$x/\delta$, and there are finitely many such vectors; by freeness distinct vectors are distinct
elements.
\end{proof}

\begin{proposition}[The lengths are everything]\label{prop:iso}
Two crystals are isomorphic --- a monoid isomorphism $\varphi$ with $\ell' \circ \varphi = \ell$ ---
if and only if their multisets of axis lengths coincide. A crystal \emph{is} its multiset of axis
lengths.
\end{proposition}

\begin{proof}
In a free commutative monoid the irreducible elements are exactly the generators, and a monoid
isomorphism carries irreducibles to irreducibles; by freeness it is determined by the induced axis
bijection. The length condition says precisely that this bijection preserves lengths, which is
possible if and only if the length multisets agree; conversely any length-preserving axis bijection
extends multiplicatively to an isomorphism.
\end{proof}

Two constructions are used throughout. The \emph{scaling} $C_\lambda = (M, \lambda\ell)$ for
$\lambda > 0$, which normalizes conventions (below, $\zeta_{C_\lambda}(s) = \zeta_C(\lambda s)$);
and the \emph{sub-crystal} on any subset $Q \subseteq P$, the free submonoid it generates with the
restricted lengths.

\begin{remark}[Attribution]\label{rem:attrib}
Definition~\ref{def:crystal} is classical. It is Beurling's system of generalized primes
\cite{Beurling, DiamondZhang} in monoid form, and it is Knopfmacher's \emph{arithmetical semigroup}
\cite{Knopfmacher} written additively ($\ell = \log$ of Knopfmacher's norm); the lattice case below
is the \emph{additive arithmetical semigroup} of Knopfmacher--Zhang \cite{KnopfmacherZhang}. The
word ``crystal'' adds only the geometric reading defended in the essay \cite{Essay}: exponent
vectors as lattice points, axes as directions, $\ell$ as the projection to a line. This note's
contribution is organization, not mathematics.
\end{remark}

\section{The standing instances}

\begin{example}[The integer crystal $C_{\Z}$]
Axes the rational primes, $\ell(p) = \log p$. Then $M \cong (\N_{\ge 1}, \times)$, $|n| = n$, and
the shadow is the set $\{\log n\}$ --- the numbers in their usual order, logarithmically spaced.
Axiom (F) holds since $\pi(x)$ is finite. All axis lengths are distinct.
\end{example}

\begin{example}[The function-field crystal $C_{\F_q[t]}$]
Axes the monic irreducibles of $\F_q[t]$, $\ell(f) = (\deg f)\log q$, so $|f| = q^{\deg f}$. Every
axis length is an integer multiple of the quantum $\log q$, and lengths are massively repeated: all
$I_n$ irreducibles of degree $n$ share one length.
\end{example}

\begin{example}[Beurling crystals]
Any sequence $1 < |p_1| \le |p_2| \le \cdots \to \infty$ of real sizes defines a crystal on
countably many axes with $\ell(p_i) = \log|p_i|$; this is exactly Beurling's construction
\cite{Beurling}, and the quantitative theory of which length regularity buys which theorems is the
subject of \cite{DiamondZhang, DMV}. The integers and the function fields are the two named points
of this continuum; the essay's temperament ladder \cite{Essay} walks between them.
\end{example}

\begin{example}[Graph crystals --- preview]
For a finite connected graph, take as axes the primitive backtrackless closed geodesics up to
rotation, with $\ell(\gamma)$ the edge length. The Euler product below becomes the Ihara zeta
function \cite{Terras}; lengths are quantized (a lattice crystal), and the Riemann Hypothesis for
the Ihara zeta is \emph{equivalent} to the graph being Ramanujan. Developed as work package WP9a.
\end{example}

\section{The invariant package}

Throughout, $C = (M, \ell)$ is a crystal. Each invariant below is a function of the length multiset
(Proposition~\ref{prop:iso}); each chapter of the essay \cite{Essay} computes one of them at
$C = C_{\Z}$.

\subsection{Zeta and partition function}

\[
\zeta_C(s) \;=\; \sum_{m \in M} e^{-s\,\ell(m)} \;=\; \sum_{m \in M} |m|^{-s}
\;=\; \prod_{p \in P} \bigl(1 - e^{-s\,\ell(p)}\bigr)^{-1},
\]
the Euler product being freeness restated. Its abscissa of convergence $\sigma_C \in [0, \infty]$ is
the crystal's \emph{density exponent}; scaling normalizes any finite nonzero $\sigma_C$ to $1$. For
$C_{\Z}$ this is Riemann's zeta; for $C_{\F_q[t]}$ it is $\zeta(s) = (1 - q^{1-s})^{-1}$, with
$I_n$ recovered by Gauss's formula. Read at real $s = \beta$, $\zeta_C(\beta)$ is the partition
function of the free bosonic gas with one mode per axis at energy $\ell(p)$ --- the primon gas
\cite{Julia}; this reading, and the phase structure of the interacting system of
Bost--Connes \cite{BostConnes}, is work package WP4.

\subsection{Layers and the tower exponent}

The \emph{dimension} grading $\Omega \colon M \to \N$ is the additive extension of $\Omega(p) = 1$.
The $k$-th \emph{layer count} of $m$ is the number $D_k(m)$ of ordered $k$-tuples of non-identity
elements with product $m$; formally
$\sum_m D_k(m)\, |m|^{-s} = (\zeta_C(s) - 1)^k$. The \emph{tower series} is the geometric sum
$G_C(s) = \sum_{k \ge 0} (\zeta_C - 1)^k = (2 - \zeta_C(s))^{-1}$, counting all ordered
factorizations of all elements.

\begin{proposition}[Tower exponent]\label{prop:tower}
Suppose $\sigma_C < \infty$ and $\lim_{s \downarrow \sigma_C} \zeta_C(s) > 2$. Then there is a
unique real $\rho_C > \sigma_C$ with $\zeta_C(\rho_C) = 2$.
\end{proposition}

\begin{proof}
On $(\sigma_C, \infty)$ the series converges to a continuous, strictly decreasing function (each
term is strictly decreasing; $M \ne \{1\}$), and $\zeta_C(s) \to 1$ as $s \to \infty$ by dominated
convergence. The intermediate value theorem gives existence; strict monotonicity gives uniqueness.
\end{proof}

For $C_{\Z}$, $\rho_C = 1.728647\ldots$ is Kalm\'ar's exponent \cite{Kalmar}, governing
$G(y) \sim Cy^{\rho}$; for $C_{\F_q[t]}$, $\zeta_C(\rho) = 2$ solves exactly to
$\rho = 1 + \log 2/\log q$ --- the integer dimension of the essay's mirror table. The Moran reading
of $\zeta_C(s) = 2$ as a similarity dimension, and the lattice/non-lattice structure of the complex
solutions, follow Lapidus \cite{LapidusBook}.

\subsection{The prime detector}

For $m \ne 1$ define the \emph{Linnik weight}
\[
\kappa(m) \;=\; \sum_{k \ge 1} \frac{(-1)^{k-1}}{k}\, D_k(m),
\]
a finite sum terminating at $k = \Omega(m)$.

\begin{proposition}[Universal detector; stated here, proved at WP2]\label{prop:linnik}
In every crystal, $\kappa(p^{\,j}) = 1/j$ on axis powers and $\kappa(m) = 0$ otherwise; the
corresponding von Mangoldt weight is $\Lambda_C(m) := \kappa(m)\,\ell(m)$, equal to $\ell(p)$ on
powers of $p$ and $0$ elsewhere.
\end{proposition}

The point of stating this at crystal generality: $D_k(p^{\,j})$ counts compositions of $j$ into $k$
positive parts, $\binom{j-1}{k-1}$, so the claim on axis powers is the identity
$\sum_k \frac{(-1)^{k-1}}{k}\binom{j-1}{k-1} = \frac1j$ (at $j = 3$:
$1 - \tfrac12\cdot 2 + \tfrac13\cdot 1 = \tfrac13$), and the vanishing off axis powers is pure
freeness --- no property of the integers enters. The formal-algebra proof (the logarithm in the
reduced incidence algebra, equivalently $\log \zeta_C$ coefficientwise) is work package WP2; at
$C_{\Z}$ it is Linnik's identity \cite{Linnik}, the engine of \cite{CompanionI}.

\subsection{Shadows and diffraction}

Two weighted combs on $[0, \infty)$:
\[
\nu_C \;=\; \sum_{m} \delta_{\ell(m)} \qquad\text{(counting comb)}, \qquad
\mu_C \;=\; \sum_{m} \Lambda_C(m)\, \delta_{\ell(m)} \qquad\text{(prime comb)}.
\]
The \emph{diffraction} of $C$ is the Fourier-transform data of $\mu_C$ in the autocorrelation sense
of mathematical diffraction theory \cite{BaakeGrimm}; making the essay's Chapter~4 exact in that
frame --- for $C_{\Z}$, the Guinand--Weil explicit formula pairing $\mu_C$ with the zeta zeros ---
is work package WP3. For lattice crystals the diffraction is periodic; for $C_{\Z}$ it is the
aperiodic spectrum conjecturally confined to the mirror line.

\subsection{Symmetry}

$\Aut(M) = \Sym(P)$: the bare crystal is maximally symmetric. The \emph{measured} automorphism
group is
\[
\Aut(C) \;=\; \{\sigma \in \Sym(P) : \ell \circ \sigma = \ell\}
\;=\; \prod_{\text{length classes}} \Sym(\text{class}),
\]
the permutations of axes within equal-length classes.

\begin{remark}[The symmetry asymmetry]\label{rem:sym}
For $C_{\Z}$ all lengths $\log p$ are distinct, so $\Aut(C_{\Z}) = 1$: \emph{measured by its sizes
alone, the integer crystal is already perfectly rigid} --- before addition is welded on. For
$C_{\F_q[t]}$, $\Aut(C) = \prod_n \Sym(I_n)$ is enormous, and it contains the arithmetic's actual
symmetries: every substitution $t \mapsto at + b$ permutes the irreducibles of each degree (the
essay's rigidity row), and the Galois action on coefficients sits inside it likewise. The essay's
thesis that the mirror is solvable \emph{because} its arithmetic has symmetries acquires a
definition-level restatement: quantized lengths leave room for symmetry ($I_n$-fold degeneracy per
level); incommensurable lengths leave none. The weld statement --- $\Aut(\N, +, \times) = 1$ ---
is logically separate and stronger; the observation here is that even the weld-free measured
crystal distinguishes the two universes.
\end{remark}

\subsection{The dichotomy}

$C$ is \emph{lattice} if $\ell(P) \subseteq \delta\Z$ for some $\delta > 0$, else
\emph{non-lattice} \cite{LapidusBook}. Lattice: $C_{\F_q[t]}$, graph crystals. Non-lattice:
$C_{\Z}$ (a rational relation among $\log p$ would exhibit a nontrivial product of primes equal to
$1$). This single bit is the essay's ``measurable difference'' row, and the Beurling continuum
parametrizes the crossing.

\section{The dictionary, restated}

\begin{center}\small
\begin{tabular}{lccc}
\hline
invariant & $C_{\Z}$ & $C_{\F_q[t]}$ & graph crystal (WP9a) \\
\hline
$\zeta_C$ & Riemann zeta & $(1 - q^{1-s})^{-1}$ & Ihara zeta \\
$\sigma_C$ & $1$ & $1$ & known \\
$\rho_C$ ($\zeta_C = 2$) & $1.7286\ldots$ (Kalm\'ar) & $1 + \log 2/\log q$ & --- \\
lattice? & no & yes ($\delta = \log q$) & yes \\
$\Aut(C)$ & $1$ & $\prod_n \Sym(I_n)$ & graph-dependent \\
detector $\Lambda_C$ & von Mangoldt & function-field $\Lambda$ & geodesic counting \\
diffraction of $\mu_C$ & zeta zeros (RH: on the mirror) & periodic comb & poles; RH $=$ Ramanujan \\
\hline
\end{tabular}
\end{center}

\noindent
The essay's two-column mirror table \cite{Essay} is this table's first two columns: one list of
invariants, two length multisets. Beurling crystals fill in the horizontal continuum.

\section{Who consumes what}

WP2 (algebraic engine): Proposition~\ref{prop:linnik} proved in the reduced incidence algebra, at
this generality. WP3 (diffraction): $\mu_C$ and its autocorrelation made rigorous; Guinand--Weil at
$C_{\Z}$. WP4 (statistical mechanics): $\zeta_C(\beta)$, $G_C(\beta)$, and $\rho_C$ as Hagedorn
point; Bost--Connes at $C_{\Z}$. WP5 (logic): decidability of the bare crystal versus the weld.
WP6 (observer ladder): correlation classes on the shadow $\nu_C$. WP8 (frontier): what extra
structure on $C_{\Z}$ replaces the mirror's $\Aut$-residue --- $\lambda$-ring data as the shadow of
the missing Frobenius. WP9a (graphs): the fourth column. The companion papers
\cite{CompanionI, CompanionII, CompanionIII} are, in this vocabulary, exact evaluations of
functionals of $\mu_{C_\Z}$ restricted to windows.

\section*{Disclosure of AI use}

This note was developed in an extensive collaboration with Claude (Anthropic; Claude Fable 5,
accessed through Claude Code, 2026). The direction of the work and all decisions about scope and
claims are the author's; the AI system drafted the text, checked the elementary proofs, and
assembled the attributions, the reason for its use being to enable a single non-specialist author
to develop and document the material to a publishable standard. The author has reviewed the note
and takes full responsibility for its content. No numerical claims are made here beyond constants
quoted from the cited companions, where they are machine-verified.

\begin{thebibliography}{99}
\bibitem{CompanionI} C. Gribble, Prime counts and prime sums from multiplication tables: a
self-similar formula system, preprint, Zenodo (2026), \texttt{doi:10.5281/zenodo.21769103}.
\bibitem{CompanionII} C. Gribble, Certified exact sums of primes over dyadic intervals via a ladder
of generalized factorials, preprint, Zenodo (2026), \texttt{doi:10.5281/zenodo.21769105}.
\bibitem{CompanionIII} C. Gribble, Recovering the primes in a dyadic interval from exact moments:
two constructive routes and an obstruction theorem, preprint, Zenodo (2026),
\texttt{doi:10.5281/zenodo.21769107}.
\bibitem{Essay} C. Gribble, The multiplication crystal: dimension, diffraction, and the symmetries
of arithmetic, manuscript (2026), in the code repository
\texttt{https://github.com/carlgribble-caa/prime-moments}.
\bibitem{Beurling} A. Beurling, Analyse de la loi asymptotique de la distribution des nombres
premiers g\'en\'eralis\'es I, \emph{Acta Math.} \textbf{68} (1937), 255--291.
\bibitem{DiamondZhang} H.~G. Diamond and W.-B. Zhang, \emph{Beurling Generalized Numbers}, Math.
Surveys Monogr. 213, Amer. Math. Soc., Providence, 2016.
\bibitem{DMV} H.~G. Diamond, H.~L. Montgomery, and U.~M.~A. Vorhauer, Beurling primes with large
oscillation, \emph{Math. Ann.} \textbf{334} (2006), 1--36.
\bibitem{Knopfmacher} J. Knopfmacher, \emph{Abstract Analytic Number Theory}, North-Holland,
Amsterdam, 1975.
\bibitem{KnopfmacherZhang} J. Knopfmacher and W.-B. Zhang, \emph{Number Theory Arising from Finite
Fields: Analytic and Probabilistic Theory}, Marcel Dekker, New York, 2001.
\bibitem{Kalmar} L. Kalm\'ar, \"Uber die mittlere Anzahl der Produktdarstellungen der Zahlen,
\emph{Acta Litt. Sci. Szeged} \textbf{5} (1931), 95--107.
\bibitem{Linnik} Yu.~V. Linnik, \emph{The Dispersion Method in Binary Additive Problems}, Amer.
Math. Soc., Providence, 1963.
\bibitem{LapidusBook} M.~L. Lapidus and M. van Frankenhuijsen, \emph{Fractal Geometry, Complex
Dimensions and Zeta Functions}, 2nd ed., Springer, New York, 2013.
\bibitem{BaakeGrimm} M. Baake and U. Grimm, \emph{Aperiodic Order. Vol. 1: A Mathematical
Invitation}, Cambridge Univ. Press, 2013.
\bibitem{Julia} B. Julia, Statistical theory of numbers, in \emph{Number Theory and Physics},
Springer Proc. Phys. 47, Springer, 1990, 276--293.
\bibitem{BostConnes} J.-B. Bost and A. Connes, Hecke algebras, type III factors and phase
transitions with spontaneous symmetry breaking in number theory, \emph{Selecta Math.} \textbf{1}
(1995), 411--457.
\bibitem{Terras} A. Terras, \emph{Zeta Functions of Graphs: A Stroll through the Garden}, Cambridge
Univ. Press, 2011.
\end{thebibliography}

\end{document}
